%% Created by Maple 2022.1, Windows 10
%% Source Worksheet: lab_sols02
%% Generated: Wed May 08 10:38:19 BST 2024
\documentclass{article}
\usepackage{amssymb}
\usepackage{graphicx}
\usepackage{hyperref}
\usepackage{listings}
\usepackage{mathtools}
\usepackage{maple}
\usepackage[utf8]{inputenc}
\usepackage[svgnames]{xcolor}
\usepackage{amsmath}
\usepackage{breqn}
\usepackage{textcomp}
\begin{document}
\lstset{basicstyle=\ttfamily,breaklines=true,columns=flexible}
\pagestyle{empty}
\DefineParaStyle{Maple Bullet Item}
\DefineParaStyle{Maple Heading 1}
\DefineParaStyle{Maple Warning}
\DefineParaStyle{Maple Heading 4}
\DefineParaStyle{Maple Heading 2}
\DefineParaStyle{Maple Heading 3}
\DefineParaStyle{Maple Dash Item}
\DefineParaStyle{Maple Error}
\DefineParaStyle{Maple Title}
\DefineParaStyle{Maple Ordered List 1}
\DefineParaStyle{Maple Text Output}
\DefineParaStyle{Maple Ordered List 2}
\DefineParaStyle{Maple Ordered List 3}
\DefineParaStyle{Maple Normal}
\DefineParaStyle{Maple Ordered List 4}
\DefineParaStyle{Maple Ordered List 5}
\DefineCharStyle{Maple 2D Output}
\DefineCharStyle{Maple 2D Input}
\DefineCharStyle{Maple Maple Input}
\DefineCharStyle{Maple 2D Math}
\DefineCharStyle{Maple Hyperlink}
\begin{center}
\begin{Maple Normal}
\underline{\textbf{Solving equations}}
\end{Maple Normal}
\end{center}
\_\_\_\_\_\_\_\_\_\_\_\_\_\_\_\_\_\_\_\_\_\_\_\_\_\_\_\_\_\_\_\_\_\_\_\_\_\_\_\_\_\_\_\_\_\_\_\_\_\_\_\_\_\_\_\_\_\_\_\_\_\_\_\_\_\_\_\_\_\_\_\_\_\_\_\_\_\_\_\_\_\textbf{Exercise 1.1}\begin{lstlisting}
> restart;
\end{lstlisting}
\begin{lstlisting}
> f := (x) -> 2*x^4-2222*x^3+224220*x^2-2222000*x+2000000;
\end{lstlisting}
\begin{lstlisting}
> solve(f(x)=0,\{x\});
\end{lstlisting}
Whenever a polynomial 
$ f (x) $ has a root 
$ x =a  $, the term 
$ x -a  $ is a factor of 
$ f (x) $.  This indicates that 
$ f (x) $ should be a multiple of the polynomial 
$ g (x)=(x -1) (x -10) (x -100) (x -1000) $.  However, the coefficient of 
$ x^{4} $ in 
$ f (x) $ is 
$ 2 $, whereas the coefficient of 
$ x^{4} $ in the expansion of 
$ g (x) $ is 
$ 1 $.  Thus, the multiplier must be 
$ 2 $, and we must have\begin{center}

$ f (x)=2 g (x) $
$ =2 (x -1) (x -10) (x -100) (x -1000) $\end{center}
Maple confirms this as follows:\begin{lstlisting}
> factor(f(x));
\end{lstlisting}
\_\_\_\_\_\_\_\_\_\_\_\_\_\_\_\_\_\_\_\_\_\_\_\_\_\_\_\_\_\_\_\_\_\_\_\_\_\_\_\_\_\_\_\_\_\_\_\_\_\_\_\_\_\_\_\_\_\_\_\_\_\_\_\_\_\_\_\_\_\_\_\_\_\_\_\_\_\_\_\_\_\textbf{Exercise 1.2}\begin{lstlisting}
> restart;
\end{lstlisting}
\begin{lstlisting}
> y := x^4-x^3-x^2-x/8+1/64;
\end{lstlisting}
\begin{lstlisting}
> _EnvExplicit := true:
\end{lstlisting}
\begin{lstlisting}
> solve(y=0,x);
\end{lstlisting}
\begin{lstlisting}
> sols := solve(y=0,\{x\});
\end{lstlisting}
In each of the four solutions, the sign attached to 
$ \sqrt{6} $ is the product of the signs attached to 
$ \sqrt{2} $ and 
$ \sqrt{3} $.  Using this observation, we see that the four solutions are (1±
$ \sqrt{2} $)(1±
$ \sqrt{3} $)/4.\begin{lstlisting}
> z := 16*x^3-24*x^2-6*x+2;
\end{lstlisting}
\begin{lstlisting}
> simplify(subs(sols[1],z),symbolic);
\end{lstlisting}
\begin{lstlisting}
> simplify(subs(sols[2],z),symbolic);
\end{lstlisting}
\begin{lstlisting}
> simplify(subs(sols[3],z),symbolic);
\end{lstlisting}
\begin{lstlisting}
> simplify(subs(sols[4],z),symbolic);
\end{lstlisting}
\begin{lstlisting}
> seq(simplify(subs(sols[i],z),symbolic),i=1..4);
\end{lstlisting}
\begin{lstlisting}
> 
\end{lstlisting}
\begin{lstlisting}
> 
\end{lstlisting}
\_\_\_\_\_\_\_\_\_\_\_\_\_\_\_\_\_\_\_\_\_\_\_\_\_\_\_\_\_\_\_\_\_\_\_\_\_\_\_\_\_\_\_\_\_\_\_\_\_\_\_\_\_\_\_\_\_\_\_\_\_\_\_\_\_\_\_\_\_\_\_\_\_\_\_\_\_\_\_\_\_\textbf{Exercise 1.3}\begin{lstlisting}
> restart;
\end{lstlisting}
\begin{lstlisting}
> _EnvExplicit := true;
\end{lstlisting}
\begin{lstlisting}
> y := x^3-3*x+1;
\end{lstlisting}
\begin{lstlisting}
> solve(y=0,\{x\});
\end{lstlisting}
\begin{lstlisting}
> sols := fsolve(y=0,\{x\});
\end{lstlisting}
\begin{lstlisting}
> plot(y,x=-2..2);
\end{lstlisting}
\begin{Maple Normal}
The only negative root is 
$ x =- 1.879 $ , which is \textcolor[RGB]{255,0,0}{\textbf{ sols[1]}}.  Maple notation for the gradient 
$ \frac{\mathit{dy}}{\mathit{dx}} $ is \textcolor[RGB]{255,0,0}{\textbf{ diff(y,x)}}.  We can thus find the gradient at the negative root as follows:
\end{Maple Normal}
\begin{lstlisting}
> subs(sols[1],diff(y,x));
\end{lstlisting}
\_\_\_\_\_\_\_\_\_\_\_\_\_\_\_\_\_\_\_\_\_\_\_\_\_\_\_\_\_\_\_\_\_\_\_\_\_\_\_\_\_\_\_\_\_\_\_\_\_\_\_\_\_\_\_\_\_\_\_\_\_\_\_\_\_\_\_\_\_\_\_\_\_\_\_\_\_\_\_\_\_\textbf{Exercise 1.4}\begin{lstlisting}
> g := (x) -> (b^2 - c^2 + (1 + c^2)*x)/(1-c^2 + c^2*x);
\end{lstlisting}
\begin{Maple Normal}
\textbf{(a)}
\end{Maple Normal}
\begin{lstlisting}
> sols := solve(g(x)=x,\{x\});
\end{lstlisting}
\begin{Maple Normal}
\textbf{(b)}
\end{Maple Normal}
The above solutions involve division by 
$ c  $, so they do not make sense when 
$ c =0 $.  We must therefore check separately what happens when 
$ c =0 $.  In that case, 
$ g (x) $ simplifies as follows:\begin{lstlisting}
> subs(c=0,g(x));
\end{lstlisting}
The equation 
$ g (x)=x  $ thus becomes 
$ x +b^{2}=x  $.  This is always true if 
$ b =0 $, and never true if 
$ b \neq 0 $.\

\textbf{(c)} Now return to the case 
$ c \neq 0 $.  The two solutions can then be written as 
$ x =1-\frac{b}{c} $ and 
$ x =1+\frac{b}{c} $.  These are the same (and both equal to 
$ 1 $) if 
$ b =0 $, but are different otherwise.  \

\

\textbf{(d)} Our conclusion is as follows:If 
$ b =0 $
$ =c  $ then 
$ g (x)=x  $ for all 
$ x  $.If 
$ b =0 $
$ \neq c  $ then the only solution to 
$ g (x)=x  $ is 
$ x =1 $.If 
$ b \neq 0 $
$ =c  $ then there are no solutions.If 
$ b \neq 0 $
$ \neq c  $ then there are precisely two different solutions, namely 
$ 1+\frac{b}{c} $ and 
$ 1-\frac{b}{c} $.\_\_\_\_\_\_\_\_\_\_\_\_\_\_\_\_\_\_\_\_\_\_\_\_\_\_\_\_\_\_\_\_\_\_\_\_\_\_\_\_\_\_\_\_\_\_\_\_\_\_\_\_\_\_\_\_\_\_\_\_\_\_\_\_\_\_\_\_\_\_\_\_\_\_\_\_\_\_\_\_\_\textbf{Exercise 2.1}\begin{lstlisting}
> f := (x) -> sin(Pi*x+exp(-x));
\end{lstlisting}
\begin{lstlisting}
> solve(f(x)=0,\{x\});
\end{lstlisting}
\begin{Maple Normal}
\textbf{(a)}
\end{Maple Normal}
\begin{lstlisting}
> plot(f(x),x=0..10);
\end{lstlisting}
The graph looks rather like a sine wave, with roots at 
$ x =1 $,2,3,4 and so on.  This makes sense, because as soon as 
$ x  $ becomes reasonably large, the term 
$ {\mathrm e}^{-x} $ becomes very small, so 
$ f (x) $ is close to 
$ \sin (\mathrm{Pi} x) $.  The positive roots of 
$ \sin (\mathrm{Pi} x) $ are exactly 
$ x =1 $,2,3,4 and so on, so we expect the roots of 
$ f (x) $ to be approximately the same.\

\

\textbf{(b)} We now find some roots more precisely:\begin{lstlisting}
> fsolve(f(x)=0,x);
\end{lstlisting}
\begin{lstlisting}
> fsolve(f(x)=0,x=2);
\end{lstlisting}
\begin{lstlisting}
> seq(fsolve(f(x)=0,x=n),n=1..10);
\end{lstlisting}
\begin{Maple Normal}
\textbf{(c)}
\end{Maple Normal}
\begin{lstlisting}
> r := (n) -> n-exp(-n)/Pi-exp(-2*n)/Pi^2;
\end{lstlisting}
\begin{lstlisting}
> seq(evalf(r(n)),n=1..10);
\end{lstlisting}
This is very close to our answer in (b).  We can calculate the differences as follows:\begin{lstlisting}
> seq(fsolve(f(x)=0,x=n)-evalf(r(n)),n=1..10);
\end{lstlisting}
\_\_\_\_\_\_\_\_\_\_\_\_\_\_\_\_\_\_\_\_\_\_\_\_\_\_\_\_\_\_\_\_\_\_\_\_\_\_\_\_\_\_\_\_\_\_\_\_\_\_\_\_\_\_\_\_\_\_\_\_\_\_\_\_\_\_\_\_\_\_\_\_\_\_\_\_\_\_\_\_\_\textbf{Exercise 3.1}\begin{Maple Normal}
\textbf{(a)}
\end{Maple Normal}
\begin{lstlisting}
> plot(tan(Pi*x)^2-3,x=-3..3,-3..3,numpoints=200);
\end{lstlisting}
\begin{Maple Normal}
\textbf{(b)}
\end{Maple Normal}
The roots are at 
$ -\frac{8}{3} $,  
$ -\frac{7}{3} $,  
$ -\frac{5}{3} $,  
$ -\frac{4}{3} $,  
$ -\frac{2}{3} $,  
$ -\frac{1}{3} $,  
$ \frac{1}{3} $,  
$ \frac{2}{3} $,  
$ \frac{4}{3} $,  
$ \frac{5}{3} $,  
$ \frac{7}{3} $,  
$ \frac{8}{3} $ and so on.  In other words, for every integer 
$ n  $, there is a root at 
$ n -\frac{1}{3} $, and another one at 
$ n +\frac{1}{3} $.\

\begin{Maple Normal}
\textbf{(c)}
\end{Maple Normal}
\begin{lstlisting}
> solve(tan(Pi*x)^2=3,\{x\});
\end{lstlisting}
\begin{Maple Normal}
\textbf{(d)}
\end{Maple Normal}
\begin{lstlisting}
> _EnvAllSolutions := true;
\end{lstlisting}
\begin{lstlisting}
> solve(tan(Pi*x)^2=3,\{x\});
\end{lstlisting}
\_\_\_\_\_\_\_\_\_\_\_\_\_\_\_\_\_\_\_\_\_\_\_\_\_\_\_\_\_\_\_\_\_\_\_\_\_\_\_\_\_\_\_\_\_\_\_\_\_\_\_\_\_\_\_\_\_\_\_\_\_\_\_\_\_\_\_\_\_\_\_\_\_\_\_\_\_\_\_\_\_\textbf{Exercise 4.1}\begin{lstlisting}
> restart;
\end{lstlisting}
\begin{lstlisting}
> eqns := \{
 x/2 + y/3 + z/4 = 1,
 x/3 + y/4 + z/5 = 2,
 x/4 + y/5 + z/6 = 3
\};
\end{lstlisting}
\begin{lstlisting}
> sol := solve(eqns,\{x,y,z\});
\end{lstlisting}
\begin{lstlisting}
> sol := solve(\{x/2+y/3+z/4=1,x/3+y/4+z/5=2,x/4+y/5+z/6=3\},\{x,y,z\});
\end{lstlisting}
\begin{lstlisting}
> subs(sol,x^2+y^2+z^2);
\end{lstlisting}
\_\_\_\_\_\_\_\_\_\_\_\_\_\_\_\_\_\_\_\_\_\_\_\_\_\_\_\_\_\_\_\_\_\_\_\_\_\_\_\_\_\_\_\_\_\_\_\_\_\_\_\_\_\_\_\_\_\_\_\_\_\_\_\_\_\_\_\_\_\_\_\_\_\_\_\_\_\_\_\_\_\textbf{Exercise 4.2}\begin{Maple Normal}
\textbf{(a)}
\end{Maple Normal}
\begin{lstlisting}
> solve(\{p+q+r=0,p+2*q+3*r=1\});
\end{lstlisting}
Note the equation 
$ p =p  $ in the solution, indicating that 
$ p  $ can take any value.  This means that there are infinitely many different solutions.\begin{Maple Normal}
\textbf{(b)}
\end{Maple Normal}
\begin{lstlisting}
> solve(\{u+v=1001,u+2*v=1002,u+3*v=1006\});
\end{lstlisting}
\begin{lstlisting}
> [solve(\{u+v=1001,u+2*v=1002,u+3*v=1006\})];
\end{lstlisting}
Maple gives an empty response, indicating that there are no solutions.  This is easy to see directly: if we subtract the first two equations we get 
$ v =1 $, whereas subtracting the second and third equations gives 
$ v =4 $, showing that the equations are not consistent.\begin{Maple Normal}
\textbf{(c)}
\end{Maple Normal}
\begin{lstlisting}
> solve(\{
  x +   y +   z = 2,
  x + 2*y + 3*z = 2,
  x + 4*y + 9*z = 2
 \},
 \{x,y,z\}
);
\end{lstlisting}
\begin{lstlisting}
> 
\end{lstlisting}
This has a unique, fully-determined solution.\end{document}
